% Generated draft: requires mathematical and source review.
\documentclass[11pt]{article}
\usepackage[margin=25mm]{geometry}
\usepackage{fontspec}
\setmainfont{Latin Modern Roman}
\usepackage{amsmath,amssymb,mathtools,unicode-math}
\setmathfont{Latin Modern Math}
\usepackage{enumitem,needspace,xurl,hyperref}
\setlength{\parindent}{0pt}
\setlength{\parskip}{5pt}
\setlength{\emergencystretch}{3em}
\newcommand{\N}{\mathbb{N}}
\newcommand{\Z}{\mathbb{Z}}
\newcommand{\Q}{\mathbb{Q}}
\newcommand{\R}{\mathbb{R}}
\newcommand{\C}{\mathbb{C}}
\newcommand{\F}{\mathbb{F}}
\newcommand{\A}{\mathbb{A}}
\newcommand{\PP}{\mathbb{P}}
\newcommand{\E}{\mathbb{E}}
\newcommand{\eps}{\varepsilon}
\newcommand{\dd}{\,\mathrm{d}}
\begin{document}
\section*{UnsolvedMath ID: 10900019}
\textbf{Source problem code:} catalogue entry
\subsection{Definitions and mathematical statement}
A geometric action is a proper cocompact action by isometries, with cubical automorphisms in the cubical case. Write
\[\begin{split}
d_{\mathrm{CAT}}(G)&=\min\{\dim X:G\text{ acts geometrically on a CAT}(0)\text{ space }X\},\\
d_{\square}(G)&=\min\{\dim X:G\text{ acts geometrically on a CAT}(0)\text{ cube complex }X\}.
\end{split}\]
The source explicitly defines the second invariant as CAT(0) cubical dimension. The first uses its corresponding CAT(0)-dimension convention, whose precise space/dimension requirements must be fixed. Whenever both finite minima exist, \(d_{\mathrm{CAT}}(G)\leq d_\square(G)\).

Walsh asks:
\begin{enumerate}
\item Does there exist a sequence \((M_i)_{i\geq1}\) of closed hyperbolic three-manifolds with
\[d_\square(\pi_1(M_i))\longrightarrow\infty?\]
Explicitly, for every integer \(N\) there is \(i_0\) such that \(d_\square(\pi_1(M_i))>N\) for all \(i\geq i_0\).
\item More generally, are there CAT(0) groups \((G_i)\) for which the difference between CAT(0) dimension and cubical dimension tends to infinity? In the finite-gap interpretation this is
\[d_\square(G_i)-d_{\mathrm{CAT}}(G_i)\longrightarrow\infty.\]
The order is the nonnegative excess of cubical dimension, as justified by the inclusion of cube complexes among CAT(0) spaces.
\end{enumerate}
The source does not state whether each cubical dimension is required finite. Setting it to \(+\infty\) for a group admitting no finite-dimensional geometric cubulation changes the substance of the sequence question, and a difference involving infinite dimensions may be undefined. The finite-gap formulation therefore records the natural nonvacuous version while flagging the missing finiteness convention.

The background contrasts this with Bridson's arbitrarily large gaps between geometric dimension and CAT(0) dimension. Geometric dimension refers to a different invariant, usually the least dimension of a classifying space \(K(G,1)\), and that comparison is not substituted for the cubical-dimension questions. The summary reports finite cubical gaps for hyperbolic three-manifold groups but gives no precise examples or references.
\subsection{Short English statement}
Can closed hyperbolic three-manifold groups require arbitrarily high-dimensional geometric cubulations? More generally, can the finite cubical dimension exceed CAT(0) dimension by arbitrarily much? The treatment of infinite dimensions needs clarification.
\subsection{Sources}
\begin{itemize}
\item Local UnsolvedMath registry, permanent ID \(10900019\), exact source code \texttt{catalogue entry}.
\item Walsh, attribution in Delp, Hoffoss and Manning, \emph{Problems In Groups, Geometry, and Three-Manifolds} (2015), subsection 3.11, PDF page 3: \url{https://arxiv.org/pdf/1512.04620#page=3}; abstract \url{https://arxiv.org/abs/1512.04620}. Registry access 2026-07-29; source-tex extraction recorded as checked against rendered PDF text.
\item Bridson's contrasting geometric-dimension examples and the summary's finite-gap progress are mentioned without precise citations.
\end{itemize}
\end{document}
